\section{Reproducibility and machine certificates}\label{sec:reproducibility}
The manuscript is assembled from one authenticated input archive. Its
primary mathematical authority is the final clean closure, which contains
the canonical upstream and local-matching components and the surviving
global-uniqueness component. The two other top-level archives are validation
and source-provenance evidence. Their audit verdicts support confidence;
the mathematical premises are the statements and proofs reproduced here,
the cited external theorems, and the specified finite certificates.

The accompanying manuscript archive contains the original input ZIP,
unaltered, under \pkg{reproducibility/}. The standard-library wrapper
\pkg{tools/replay_frozen.py} offers the command
\begin{center}
\texttt{python tools/replay\_frozen.py --full}
\end{center}
from the manuscript directory. It checks the manuscript inventory, the
original archive hash and its input manifest, extracts into a fresh temporary
directory, and runs the frozen component replayers without editing their
source. The quick command, omitting \pkg{--full}, checks the final integration
and its mutation tests. Full mode additionally regenerates R0's finite
identities, R1's exact witnesses, the local-tail symbolic certificate, and
the global matching and root certificates. No network access is required.
The exact command scope and environment are recorded in
\pkg{REPRODUCIBILITY.md}.

The fresh assembly checks passed the final integration replay, including
12 distinct nested archive identities, 590 nested manifest entries, the
typed identification chain, strict endpoint signs, and all 13 rational
Newton intersections. The component replays for R0, R1, local matching,
and global uniqueness also passed. R1 rejects 13 mathematical mutations.
The integration mutation checker rejects missing interval premises, changed
matching objects, branch substitution, false endpoint signs, invalid Newton
centers and other specified corruptions. The scope is finite rejection
testing, not a proof that every possible implementation defect is absent.
Fresh logs and a receipt are included in the manuscript archive.

The arithmetic layers have different contracts. R0 and R1 use arbitrary
precision integers and rational arithmetic. The local fixed-point and
characteristic identities use exact symbolic operations, with frozen
SymPy 1.14.0 and its bundled dependency. Global matching uses a small
project-local real/complex interval and first-derivative implementation,
with directed Decimal contexts at precision 90. A square root is widened
from its correctly rounded nearest value and accepted only after the
squared endpoints enclose the argument. Complex rectangles certify the
chosen square-root branch and nonzero division denominators. Analytic
remainder bounds are added before propagation; they are not inferred from
rounding alone. The final enclosure arithmetic is independently checked
using rational numbers.

The executable trust base includes the checker sources, Python and its
integer, rational and Decimal implementations, parsing and filesystem
behavior, the operating system and hardware. The stable-manifold,
compactness, contraction and variational arguments are ordinary analytic
proofs rather than executable proof-assistant terms. Hashes establish
custody and exact version identity, not mathematical correctness.
Detailed identities and certificate paths are in
Appendices~\ref{app:bounds}--\ref{app:crosswalk}; full archive hashes are
in Table~\ref{tab:hashes}.

The surviving global-uniqueness source has archive digest beginning
\pkg{106adfdeb56e}. The rejected alternative branch is not used.
The distinction is enforced by the frozen integration's closed archive
allowlist and exact source pins. Replaying a branch with a numerically
matching answer would not meet this contract.
